% eksperymentalne tłumaczenie na włoski

%

% PRZECZYTANO: https://en.wikipedia.org/wiki/Pons_asinorum

\index{pons asinorum|(}

\pol{\section{Pons asinorum}}
\ita{\section{Pons asinorum}}

\pol{,,\emph{Pons asinorum}'', czyli most osłów, to tradycyjna nazwa twierdzenia (I.5), że kąty przy podstawie trójkąta równoramiennego są równe.}
\ita{\emph{Pons asinorum}, cioè ``ponte degli asini'', è il nome tradizionale del teorema (I.5) secondo cui gli angoli alla base di un triangolo isoscele sono uguali.}
\pol{Ci, którzy nie są w stanie samodzielnie przeprowadzić jego dedukcyjnego dowodu opartego na własnościach trójkątów przystających, nie mogą przekroczyć mostu i studiować dalej geometrii.}
\ita{Chi non è in grado di seguirne autonomamente la dimostrazione deduttiva basata sulle proprietà dei triangoli congruenti non può attraversare il ponte e proseguire nello studio della geometria.}
\pol{Bardziej przyziemnie Coxeter \cite[s. 22-24]{coxeter_1967} zauważa, że rysunek wykonany przez Euklidesa przypomni most.}
\ita{Più concretamente, Coxeter \cite[pp. 22-24]{coxeter_1967} osserva che la figura disegnata da Euclide ricorda un ponte.}
\pol{Wśród konsekwencji wymienia wyniki z~Elementów: (III.3), (III.20), (III.21), (III.22), (III.32), (VI.2), (VI.4), a potem (III.35), (III.36), (VI.19), co prowadzi do dowodu twierdzenia Pitagorasa, czyli (I.47).}
\ita{Tra le conseguenze cita i risultati degli \emph{Elementi}: (III.3), (III.20), (III.21), (III.22), (III.32), (VI.2), (VI.4), e successivamente (III.35), (III.36), (VI.19), che conducono alla dimostrazione del teorema di Pitagora, ossia (I.47).}
\pol{\index{twierdzenie!Pitagorasa}}%
\ita{\index{teorema!di Pitagora}}%

\begin{figure}[H] \centering
...
\end{figure}

\pol{Pierwsze dowody tego faktu podadzą jeszcze Euklides, komentujący jego prace Proklos, a także (dużo krócej\footnote{Pappus zauważa, że trójkąt $\triangle ABC$ przystaje do siebie $\triangle ACB$, więc stosowne kąty przy podstawie też są przystajace.}) Pappus z Aleksandrii.}
\ita{Le prime dimostrazioni di questo fatto furono fornite da Euclide, dal suo commentatore Proclo e anche (molto più brevemente\footnote{Pappo osserva che il triangolo $\triangle ABC$ è congruente a $\triangle ACB$, per cui anche i corrispondenti angoli alla base sono congruenti.}) da Pappo di Alessandria.}
\index[persons]{Proklos zwany Diadochem}%
\index[persons]{Pappus z Aleksandrii}%
\pol{Przyszłość przyniesie jeszcze jedno uzasadnienie, zaczynające się od wykreślenia dwusiecznej z kąta przy wierzchołku.}
\ita{In seguito comparirà un'altra dimostrazione, che inizia tracciando la bisettrice dell'angolo al vertice.}
\pol{\index{dwusieczna}}%
\ita{\index{bisettrice}}%
\pol{Euklides nie zrobi tego przede wszystkim ze względu na kolejność wykładanego materiału: dwusieczna pojawi się cztery tezy później, a nie można korzystać z wyników, których prawdziwości dopiero się pokaże.}
\ita{Euclide non procede in questo modo soprattutto per ragioni di ordine espositivo: la bisettrice compare quattro proposizioni più tardi e non si possono usare risultati che non sono ancora stati dimostrati.}

\pol{O pons asinorum nie wspomina żaden szkolny podręcznik geometrii \texttt{:(}}
\ita{Nessun manuale scolastico di geometria menziona il \emph{pons asinorum} \texttt{:(}}
\pol{Pojawia się u Bogdańskiej, Neugebauera \cite[s. 9]{neugebauer_2018}.}
\ita{Compare invece in Bogdańska e Neugebauer \cite[p. 9]{neugebauer_2018}.}

\index{pons asinorum|)}

\pol{\index{most osłów|see{pons asinorum}}}
\ita{\index{ponte degli asini|see{pons asinorum}}}

%